% 第一章 绪论
% 说明：本文引用采用 \cite{} 形式，文献 key 可在 bib 文件中统一维护。
% 若学校模板使用数字编号制，编译后会自动显示为 [1]、[2] 等编号。

\chapter{绪论}

\section{研究背景与意义}

随着移动互联网、5G 通信、云计算和智能终端设备的快速发展，在线视频已经成为互联网内容传播和用户娱乐消费的重要形式。以短视频、长视频、直播回放和用户自制内容为代表的视频内容不断增长，视频平台逐渐从单一的视频播放工具发展为集内容创作、社交互动、搜索推荐和内容治理于一体的综合型平台。用户不仅是视频内容的观看者，也是视频内容的生产者和传播者，评论、弹幕、点赞、收藏、投币、关注等交互方式进一步增强了平台的社区属性。

在线视频分享平台在带来丰富内容生态的同时，也对系统架构提出了更高要求。一方面，视频文件通常体积较大，上传、存储、转码、切片、播放等流程涉及大量计算和 I/O 操作，需要系统具备较好的处理能力和扩展能力；另一方面，平台需要支持用户登录、视频投稿、视频检索、互动行为、后台管理等多类业务，传统单体架构在功能扩展、模块维护和服务部署方面容易出现耦合度高、扩展困难、故障影响范围大等问题。微服务架构强调将复杂系统拆分为围绕业务能力构建的小型服务，各服务可以独立开发、部署和扩展，因此适合功能模块较多、业务流程较长的在线视频平台系统\cite{lewis2014microservices,newman2021building,fielding2000architectural,evans2003ddd,richardson2018microservices}。

此外，随着用户生成内容数量的增加，平台内容安全问题也越来越突出。视频内容可能包含暴力、低俗、敏感言论或其他违规信息，仅依靠人工审核难以满足效率和实时性的要求。人工智能技术的发展，特别是语音识别、声音事件检测和多模态大模型的发展，为视频内容审核提供了新的技术方案。Whisper 可以用于大规模弱监督语音识别，AudioSet 和 YAMNet 为声音事件识别提供了数据与模型基础，Qwen-VL 系列模型则体现了多模态理解在图像、文字和语义判断中的应用价值\cite{radford2023whisper,gemmeke2017audioset,hershey2017cnn,tensorflow2024yamnet,bai2023qwen}。通过将 AI 审核服务与业务系统结合，可以在视频发布前对画面、语音文本和声音事件进行综合分析，给出风险等级和审核建议，再由管理员进行人工复核，从而形成“AI 初审 + 人工终审”的协同审核流程。

本课题设计并实现一个基于微服务架构的在线视频分享平台，系统面向普通用户、内容创作者和平台管理员，提供视频浏览、搜索、投稿、转码播放、互动评论、后台审核和 AI 智能审核等功能。课题的意义主要体现在三个方面：第一，通过 Spring Boot、Spring Cloud、Spring Cloud Gateway、Nacos、OpenFeign 等技术完成在线视频平台的微服务化设计，提高系统的可维护性和扩展性\cite{springframeworkdocs,springbootdocs,springclouddocs,springgatewaydocs,springopenfeign,nacosdocs}；第二，通过 FFmpeg、HLS、Redis、Elasticsearch、RabbitMQ 等技术实现视频处理、搜索和异步任务协同，提升系统工程完整性\cite{ffmpegdocs,rfc8216hls,redisdocs,elasticsearchdocs,rabbitmqdocs}；第三，通过独立 Python AI 审核服务实现多模态内容审核，为视频平台内容安全治理提供可行的技术实践\cite{radford2023whisper,gemmeke2017audioset,hershey2017cnn,tensorflow2024yamnet,bai2023qwen,fastapidocs,uvicorndocs,pikadocs}。

\section{国内外研究现状}

国外在线视频平台起步较早，YouTube、Netflix 等平台在视频分发、流媒体播放、内容推荐和系统架构方面积累了大量实践经验。其中，Netflix 是微服务架构较早的实践者之一，微服务架构相关研究和实践强调通过服务拆分、注册发现、网关路由、弹性伸缩和自动化部署等技术提高系统的可用性和扩展能力\cite{lewis2014microservices,newman2021building,richardson2018microservices}。随着云原生技术的发展，国外视频平台普遍采用微服务、容器化部署、分布式存储和内容分发网络等方式支撑大规模用户访问。与此同时，HLS 等自适应流媒体技术逐渐成熟，FFmpeg 等开源视频处理工具被广泛应用于视频转码、切片和格式转换等场景\cite{ffmpegdocs,rfc8216hls}。

国内在线视频平台和短视频平台发展迅速，哔哩哔哩、抖音、快手等平台在用户创作、弹幕互动、视频推荐和社区运营方面形成了较成熟的业务模式。国内企业在后端架构上也广泛采用微服务和分布式技术，常见技术包括 Spring Cloud Alibaba、Nacos、Redis、消息队列和搜索引擎等。这些技术能够较好地解决业务模块复杂、访问量大、数据读写频繁和服务扩展困难等问题。Spring Boot、Spring Cloud、Spring Cloud Gateway、Nacos、Redis、Elasticsearch 和 RabbitMQ 等官方文档也说明了这些技术在服务构建、路由治理、缓存、搜索和消息通信方面的工程价值\cite{springbootdocs,springclouddocs,springgatewaydocs,nacosdocs,redisdocs,elasticsearchdocs,rabbitmqdocs}。

在内容审核方面，国内外平台均逐渐从单纯依赖人工审核转向“机器审核 + 人工复核”的模式。传统审核方式主要依靠关键词匹配、图像分类模型或人工规则，能够处理部分简单违规内容，但面对复杂语义、音视频混合内容和隐蔽表达时存在局限。近年来，多模态人工智能技术快速发展，视频审核逐渐从单一的图像或文本检测扩展到画面、语音、文本和声音事件的综合分析。语音识别模型可以将视频中的语音转换为文本，声音事件检测模型可以识别异常声音，多模态大模型可以结合画面和文本进行语义判断，从而提升审核结果的准确性和解释性\cite{radford2023whisper,gemmeke2017audioset,hershey2017cnn,tensorflow2024yamnet,bai2023qwen}。

总体来看，当前在线视频平台相关技术已经较为成熟，但对于毕业设计项目而言，如何将微服务架构、视频处理、用户互动、搜索功能和 AI 审核服务整合为一个完整可运行的平台，仍具有较强的综合实践意义。本课题在现有技术基础上，结合 Spring Cloud Alibaba 微服务体系和 Python AI 审核服务，构建一个具备核心业务闭环的在线视频分享平台。

\section{研究内容与目标}

本课题围绕“基于微服务架构的在线视频分享平台设计与实现”展开，主要研究如何将在线视频分享平台中的用户管理、视频投稿、资源处理、互动行为、后台审核和 AI 内容审核等功能进行合理拆分与协同实现。系统整体采用前后端分离架构，前端包括普通用户端和后台管理端，后端采用 Spring Cloud Alibaba 微服务架构，AI 审核部分采用独立 Python 服务实现。前后端分离和微服务拆分能够降低模块耦合度，便于不同模块独立开发和后续扩展\cite{lewis2014microservices,newman2021building,richardson2018microservices,vueguide,viteguide,vuerouterdocs,piniadocs,axiosdocs}。

本课题的主要研究内容如下：

\begin{enumerate}
  \item 完成在线视频分享平台的需求分析与总体架构设计。根据平台业务特点，将系统划分为网关服务、主站业务服务、后台管理服务、资源文件服务、互动服务、公共模块和基础模块等部分，并通过 Nacos、Spring Cloud Gateway、OpenFeign 等技术实现服务注册、路由转发和服务间调用\cite{springclouddocs,springgatewaydocs,springopenfeign,nacosdocs}。
  
  \item 实现视频投稿、上传、转码和播放链路。用户可以在创作中心上传视频和封面，填写视频标题、分类、标签、简介等投稿信息。系统通过分片上传降低大文件上传失败的影响，通过资源服务对视频进行合并、转码和 HLS 切片处理，并在转码完成后回写视频文件状态，为后续审核和播放提供基础。该部分主要结合 FFmpeg 和 HLS 实现视频转码与流式播放\cite{ffmpegdocs,rfc8216hls,whatwgvideo,mdnvideoaudio}。
  
  \item 实现视频浏览、搜索和互动功能。用户可以在前台浏览视频列表，按照分类或关键词搜索视频，进入视频详情页进行播放，并进行评论、弹幕、点赞、收藏、投币、关注等操作。系统使用 Elasticsearch 支撑视频检索，使用 Redis 缓存部分热点数据和临时状态，提高访问效率\cite{redisdocs,elasticsearchdocs}。
  
  \item 实现后台管理与人机协同审核流程。管理员可以在后台进行分类管理、视频审核和资源查看。视频完成转码后先进入待审核状态，AI 服务对视频进行自动审核并给出风险建议，管理员结合 AI 审核结果进行人工通过或驳回，审核通过后视频正式发布。该流程将消息队列、AI 模型和人工审核结合起来，体现了异步任务处理和人机协同的系统设计思想\cite{rabbitmqdocs,radford2023whisper,gemmeke2017audioset,hershey2017cnn,tensorflow2024yamnet,bai2023qwen}。
  
  \item 实现独立的 AI 多模态审核服务。AI 服务通过 RabbitMQ 消费审核请求，下载待审核视频后执行音频提取、语音分段、声音事件检测、语音转写、关键帧抽取和多模态模型判断，并将审核建议、风险等级、风险标签、命中片段和关键帧截图等结果返回给 Java 后端\cite{radford2023whisper,gemmeke2017audioset,hershey2017cnn,tensorflow2024yamnet,bai2023qwen,silerovad,fastapidocs,uvicorndocs,pikadocs}。
\end{enumerate}

本课题的目标是完成一个功能相对完整、架构清晰、流程闭环的在线视频分享平台原型。系统应能够展示微服务架构在复杂业务拆分中的应用，能够完成视频从上传到审核再到发布的核心流程，并能够通过 AI 审核服务提高视频内容审核的效率和辅助决策能力。

\section{论文结构安排}

本文围绕在线视频分享平台的设计与实现展开，全文共分为七章，各章内容安排如下：

第一章为绪论，主要介绍课题的研究背景与意义，分析国内外在线视频平台、微服务架构和 AI 内容审核的发展现状，明确本文的研究内容、研究目标和论文结构。

第二章为相关技术介绍，主要介绍系统开发中使用的关键技术，包括微服务架构与服务治理技术、前端开发技术、数据存储与缓存技术、视频处理技术、消息队列技术和 AI 多模态审核技术。

第三章为系统需求分析，主要从用户端、创作中心、管理端和 AI 审核等方面分析系统功能需求，并对系统的可扩展性、可维护性、内容安全、性能和可行性进行分析。

第四章为系统总体设计，主要介绍系统整体架构、微服务模块划分、服务调用关系、数据库设计、缓存与搜索设计、消息通信设计以及核心业务流程设计。

第五章为系统详细设计与实现，主要对用户端功能、视频上传转码播放、互动搜索、后台管理、AI 审核任务编排、多模态审核服务和人机协同审核等核心模块的设计与实现进行说明。

第六章为系统测试与结果分析，主要介绍测试环境和测试方案，对系统核心功能、视频处理流程、AI 审核流程和异常场景进行测试，并对测试结果进行分析。

第七章为总结与展望，主要总结本文完成的工作，分析项目的创新点与不足，并对后续系统优化方向进行展望。
